\section{Reasons for skepticism and responses}
\label{sec:reasons_for_skepticism}

We have made a case that an ambitious mathematical theory of deep learning is possible and that developing this theory is a worthwhile endeavor.
This is far from a universal view, and so we now address common counterarguments that a theory of deep learning is either not possible or not a goal worthy of our effort.

\paragraph{Competent researchers have been trying to develop a theory of deep learning for decades, and we don't have one.
Surely if there was a theory, we would have already found it.}
It is true that machine learning theory is a field with a long history, and certain avenues for developing theory have been thoroughly explored.
Why should now be different?

There are several reasons for optimism.
First, the practical success of deep learning is comparatively recent, and we have a wealth of new empirical systems to study and mine for explainable phenomena.
Some of these phenomena, like the apparent convergence to universal representations discussed in \cref{sec:reason-universal-behavior}, were only revealed by the last few years of model scaling.
These developments have turned the search for a theory of deep learning from a mathematics into an empirical science (and one with no lack of interesting things to measure).
We now have much better means to ask questions and check our answers in a tight feedback loop.

Second, the field is much bigger: empirical successes have attracted researchers from physics, mathematics, neuroscience, and other adjacent fields, and so we have more and more diverse minds on the case. Third, it is worth noting that the development of major sciences has usually taken at least several decades, so we should not be too discouraged that we do not yet have all the answers.


\paragraph{The objects currently understood from theory are very primitive compared to e.g. LLMs. Surely first-principles understanding of large models is too heavy a lift.}
Indeed, we expect that building up to LLMs will be a heavy lift and take considerable time.
The near-term hope is instead that some understanding of the basic building blocks of deep learning will prove useful even without a constructive theory that explains the whole model.
We can see this happening already in isolated pockets, including empirical scaling laws (\cref{sec:reason-measurable}), mathematical prescriptions for hyperparameter scaling (\cref{sec:reason-hps}), neural-tangent-kernel-based methods for data attribution \citep{park:2023-trak}, and theoretically motivated optimizers \citep{gupta2018shampoo,jordan2024muon}.
These ``local theories'' of small pieces of the deep learning stack are useful for hyperparameter scaling in large models, even though they are in no way comprehensive theories of the model!
One might hope for similarly useful ``local theories'' that treat subjects like training instabilities, dataset selection and attribution, or the effect of normalization layers.

It is also important to stress that the identification of the right \textit{basic} objects in a field of science often makes it possible to ask \textit{applied} questions in a more sensible way.
Consider, for example, how the understanding that all matter is made of atoms underlies virtually all other basic science, and how knowledge of electromagnetism permits optical and radiological tools in countless applied disciplines.
As discussed in \cref{subsec:learning_mech_and_mechinterp}, we hope that learning mechanics can offer tools that adjacent fields such as mechanistic interpretability can apply to better carry out their work.
In this way, rigorous work on primitive objects can aid the applied science of large models even without a rigorous theory that builds all the way up.





\paragraph{What matters is a model's high-level behavior. Microscopic theories are too zoomed in to see this.}
Models' high-level behavior is indeed important.
How does this fit in with the lower-level sciences of deep learning?
We argue that deep learning may be studied at the level of physics, biology, or \textit{psychology,} with this last including the study of the model's capabilities, personality~\citep{betley2026training}, and goals.
It seems likely that study at all levels will be necessary.
Learning mechanics (the physics of deep learning) is the farthest from model psychology, with mechanistic interpretability (the biology) lying in the middle and connecting the two.%
\footnote{We note that these three levels of study of deep learning are roughly analogous to Marr's levels of analysis of a computational system: \textit{the physical implementation of the computation}, \textit{how the computation is performed algorithmically}, and \textit{what is being computed} \citep{marr2010vision}.}


\paragraph{We don't need a theory of deep learning, we need a theory of data.}
We think we need both: we need a theory of the structure in data and a theory of how a parameterized model learns it.
We touch on the necessity of developing a useful theory of data in \cref{sec:reason-universal-behavior,openquestion:theory_for_real_data}.
These are both part of the project of developing a mechanics of learning.

\paragraph{AI will understand itself before we do. Why try to build theory?}
This is a present concern for human intellectual endeavors across the board.
Our response here has three parts.
First, theory is \textit{already} useful, and will continue to be more impactful as it develops, so this scientific work is likely to make a near-term impact.
Second, it seems unlikely that AI working in isolation will suddenly and separately ``solve deep learning theory.''
It seems more likely that breakthrough progress in a transitory period will come from human scientists using or working with AI, and expert humans will remain in the loop.
Third, if one's goal is AI safety, some human oversight of AI systems will be necessary (unless one trusts the AIs to fully police themselves), and having a human-parseable theory of deep learning gives us a foot in the door.
